\section{Projects}
% \subsection{Image Generation using Multi-marginal Optimal Transport}{University of Toronto}{2022 - 2023}
% \begin{itemize}[leftmargin=*]
%     \item Developed an advanced image generation framework leveraging multi-marginal optimal transport theories, utilizing PyTorch. This project emphasizes my capability to translate complex theoretical models into practical computational algorithms.
%     \item Analyzed the CelebA dataset, comprising over 200,000 images, applying sophisticated numerical optimization techniques to improve generative model performance.
% \end{itemize}
\subsection{Generative Modeling for scRNA-seq Data}{University of Toronto}{2021 - 2023}
\begin{itemize}[leftmargin=*]
    \item Developed population dynamics models using stochastic differential equations (SDEs) and game theory to simulate cell trajectories.
    \item Published in PLOS Computational Biology, demonstrating novel methods for modeling high-dimensional single-cell data. (\href{https://github.com/Diebrate/population_model}{GitHub} and \href{https://journals.plos.org/ploscompbiol/article?id=10.1371/journal.pcbi.1012015}{paper})
\end{itemize}

\subsection{Anomaly Detection via Optimal Transport for Time Series Data}{University of Toronto}{2021 - 2023}
\begin{itemize}[leftmargin=*]
    \item Designed statistical methods for anomaly detection, with applications in biological time series, including metabolic and cellular processes.
    \item Applied optimal transport techniques related to WGAN to enhance accuracy and scalability of anomaly detection models. (\href{https://github.com/Diebrate/scRNA_study}{GitHub} and \href{hhttps://arxiv.org/abs/2501.03501}{paper}).
\end{itemize}
